\section{Time-Series Cross-Validation}\label{sec:cv}
\subsection{Protocol and aggregation}
Notebook 04 keeps the default estimators and h24 feature set D fixed while replacing the single November assessment with five expanding-window folds. This separates the effect of broader temporal evaluation from the parameter changes introduced later. Each validation window has 720 hours; training sizes increase from \num{4225} to \num{7105} hours. The exact dates appear in \cref{tab:folds}, and December is excluded. RF-P1 is not included in this stage's saved comparison.

For fold RMSEs $R_0,\ldots,R_4$, the reported summaries are
\begin{equation}
\overline R=\frac{1}{5}\sum_{k=0}^{4}R_k,\qquad
s_R=\sqrt{\frac{1}{4}\sum_{k=0}^{4}(R_k-\overline R)^2}.
\label{eq:cv_summary}
\end{equation}
The same mean/sample-SD convention applies to the other fold metrics. These are summaries of temporal variation under increasing training sizes. They are not confidence intervals, and the overlapping training windows do not constitute independent replications.

\subsection{Average performance and stability}
\Cref{tab:cv} changes the conclusion suggested by November alone. LR has the lowest default mean RMSE, \num{2792.3} kW, with SD \num{623.3} kW. RF and XGB follow at \num{3069.5} and \num{3264.8} kW, with much larger SDs of \num{1488.8} and \num{1815.8} kW. Naive-24h reaches \num{3392.4} kW with SD \num{421.5} kW. LR is the most stable default ML model by this dispersion measure, while Naive-24h has still lower dispersion overall.
\begin{table}[H]\centering\small
\caption{Default-model development CV, h24 (T05). SD is the sample standard deviation over five folds, not a confidence interval. Errors in kW; MAPE in percent.}\label{tab:cv}
\setlength{\tabcolsep}{4pt}
\begin{tabular}{lrrrr}
\toprule
Model & RMSE mean $\pm$ SD & MAE & MAPE & R$^2$ \\
\midrule
LR & \num{2792.3} $\pm$ \num{623.3} & \num{2101.0} & \num{2.95} & \num{0.9651} \\
RF & \num{3069.5} $\pm$ \num{1488.8} & \num{2168.6} & \num{2.82} & \num{0.9582} \\
XGB & \num{3264.8} $\pm$ \num{1815.8} & \num{2328.0} & \num{2.99} & \num{0.9510} \\
Naive-24h & \num{3392.4} $\pm$ \num{421.5} & \num{2371.7} & \num{3.28} & \num{0.9497} \\
SVR & \num{3968.7} $\pm$ \num{2891.5} & \num{2753.4} & \num{3.44} & \num{0.9224} \\
DT & \num{4048.2} $\pm$ \num{1410.4} & \num{2915.3} & \num{3.87} & \num{0.9294} \\
Naive-168h & \num{4892.6} $\pm$ \num{2297.1} & \num{3655.4} & \num{4.92} & \num{0.8887} \\
\bottomrule\end{tabular}
\end{table}


Default SVR and DT have mean RMSEs of \num{3968.7} and \num{4048.2} kW, exceeding Naive-24h. This does not contradict SVR's strong November result: November is only the final fold. Nor does LR's first-place RMSE ranking extend to every metric. RF has the lowest mean MAPE, \num{2.82}\%, compared with LR's \num{2.95}\%. The implemented selection objective remains mean RMSE rather than a retrospective choice of the most favorable metric; the repository does not establish prospective preregistration.

\Cref{tab:cv_folds,fig:cv_folds} expose the fold structure behind these averages. The last fold reproduces the November results for all eight shared models and baselines within the stored numerical tolerance. Among the five ML families, LR rises from fourth in November to first by mean CV RMSE; XGB moves from first to third and SVR from second to fourth. Across all eight entries, SVR is fifth because Naive-24h also ranks ahead of it. Stating the comparison population avoids confusing ML-only and combined rankings.
\begin{table}[tbp]\centering\small
\caption{Default h24 RMSE by temporal fold (T05), in kW. Models are ordered by mean fold RMSE; fold boundaries are given in Table~\ref{tab:folds}.}\label{tab:cv_folds}
\setlength{\tabcolsep}{4pt}
\begin{tabular}{lrrrrr}
\toprule
Model & Fold 0 & Fold 1 & Fold 2 & Fold 3 & Fold 4 \\
\midrule
LR & \num{3070.8} & \num{3307.8} & \num{3323.4} & \num{1970.8} & \num{2288.7} \\
RF & \num{5397.2} & \num{3456.3} & \num{2780.0} & \num{1486.8} & \num{2227.5} \\
XGB & \num{6229.6} & \num{3433.3} & \num{2988.7} & \num{1510.9} & \num{2161.4} \\
Naive-24h & \num{4024.4} & \num{3562.6} & \num{3328.4} & \num{2970.0} & \num{3076.4} \\
SVR & \num{8774.7} & \num{4352.6} & \num{3040.3} & \num{1480.6} & \num{2195.2} \\
DT & \num{6110.5} & \num{4663.4} & \num{3906.2} & \num{2568.3} & \num{2992.4} \\
Naive-168h & \num{6224.7} & \num{7205.0} & \num{6062.4} & \num{1768.2} & \num{3202.9} \\
MeanProfile & \num{21060.3} & \num{14973.6} & \num{6654.0} & \num{7967.9} & \num{9369.6} \\
\bottomrule\end{tabular}
\end{table}

\begin{figure}[tbp]\centering
\includegraphics[width=\linewidth]{F12_rmse_per_fold.png}
\caption{Default h24 RMSE across expanding validation folds (F12). Fold 0 covers 4 July--2 August and reveals substantial sensitivity of several nonlinear estimators to the summer period. MeanProfile is omitted from the graphic but retained in the numerical tables.}\label{fig:cv_folds}
\end{figure}

\subsection{Summer distribution shift and model behavior}
The first fold is particularly difficult for default RF, XGB, SVR and DT, with RMSEs approximately \num{5397}, \num{6230}, \num{8775} and \num{6111} kW. LR and Naive-24h reach approximately \num{3071} and \num{4024} kW. The audit's hourly-data check finds that the largest training load is \num{110710.5} kW, while the largest validation load is \num{133129.8} kW; \num{12.22}\% of validation hours exceed the training maximum. This percentage is the fraction of observations above that maximum, not the percentage difference between the two maxima.

Standard regression-tree predictions are averages of training responses within leaves; an RF averages such predictions. Their absolute predictions therefore cannot exceed the training response range. This restriction makes unusually high validation loads challenging for those models. The same strict range bound does not apply generally to additive XGBoost or RBF SVR. Their summer errors are compatible with difficulty adapting to a shifted regime, but the saved experiments do not prove a unique failure mechanism for each estimator. The residual-target formulation in \cref{sec:tuning} tests a practical response to this difficulty.

Training and validation errors reinforce this qualification. The default tree interpolates training data, while RF and XGB have mean training RMSEs around \num{797.2} and \num{697.8} kW, far below their mean CV errors. Both overfitting and temporal distribution change can contribute to such gaps. A validation/training ratio cannot disentangle them by itself.

\subsection{Implications for subsequent selection}
The temporal folds demonstrate that November alone would favor a different default model. They also show that a recent-load baseline remains competitive when nonlinear defaults encounter a new regime. This motivates tuning regularization and comparing absolute with lag-relative targets using the common folds. However, these same development folds have already informed feature and modeling decisions. Their later reuse in search introduces selection uncertainty that fold SD does not quantify. The separate December evaluation must retain the CV selection rule and cannot convert these development scores into unbiased estimates of future performance.
\evidence{Notebook 04 cells 6--19; \code{T05_level2_cv.csv}; \code{T05_level2_summary.csv}; \code{T02_cv_folds.csv}; F12; \code{report/evidence/numerical_checks.json} for the fraction above the fold-0 training maximum.}
